% !TEX root = paper.tex
\doclearpage
\section{String Arithmetization}
\label{sec:string-arith}
Our goal int this paper is to design protocols for \emph{white-box} string operations and the first barrier that we encounter is that $\truthtable$ arithmetizes a string column as a column of hashes, which are inherently \emph{black-box}. A hash exposes nothing about the string it stands for, so no protocol can reason about the string's contents. Therefore, the first step is to design a new arithmetization that exposes the characters of the string columns.

\ifsp\begin{figure*}[tp]\else\begin{figure}[h!]\fi
\centering
% Text is set at the body font size: the picture is not rescaled; instead
% the grid unit \u is chosen so the figure fits \textwidth.
\ifsp\def\u{0.72}\else\def\u{0.67}\fi
% one colour per string; text uses the full colour, cells a light tint
\colorlet{applecol}{red!75!black}
\colorlet{orangecol}{orange!90!black}
\colorlet{cherrycol}{purple!85!black}
\colorlet{bananacol}{olive}
\begin{tikzpicture}[
    x=\u cm, y=\u cm,
    cell/.style={draw, minimum size=\u*6mm, inner sep=0pt},
    rowlab/.style={font=\itshape, anchor=east},
    hd/.style={font=\bfseries},
    boxlab/.style={font=\itshape, align=center, anchor=north},
    idx/.style={font=\scriptsize, gray}
  ]

  %==================== Panel 1: str-level columns ====================
  \node[hd, anchor=west] at (-0.9,5.2) {$\mathsf{s}$};
  \node[hd] at (2.1,5.2) {$\mathsf{a}$};
  \node[hd] at (2.7,5.2) {$\mathsf{h}$};
  \node[hd] at (4.4,5.2) {$\mathsf{l}$};
  \node[hd] at (5.0,5.2) {$\mathsf{rk}$};

  % i / string / length / a_h / rank ; row 2 ("cherry") is inactive.
  % Ranks are dense ranks in lexicographic order: apple < banana < cherry < orange.
  \foreach \i/\s/\ll/\ah/\rk/\col in {0/apple/5/1/0/applecol, 1/orange/6/1/3/orangecol, 2/cherry/6/0/2/cherrycol, 3/banana/6/1/1/bananacol}{
    \pgfmathsetmacro{\yy}{4.6 - \i*0.6}
    \node[font=\ttfamily, anchor=west, text=\col] at (-0.9,\yy) {\s};
    \node[cell, fill=\col!15] at (2.1,\yy) {\ah};
    \node[cell, fill=\col!15] at (2.7,\yy) {$h_{\i}$};
    \node[cell, fill=\col!15] at (4.4,\yy) {\ll};
    \node[cell, fill=\col!15] at (5.0,\yy) {\rk};
  }

  % dotted box bundling the main column (its activator and hash data)
  \draw[black, thick, densely dotted, rounded corners=2pt] (1.75,5.7) rectangle (3.05,2.2);
  \node[boxlab, anchor=north east, align=center] at (3.3,2.05) {main\\columns};

  % dotted box around the str-level auxiliary columns (length, rank)
  \draw[black, thick, densely dotted, rounded corners=2pt] (4.05,5.7) rectangle (5.35,2.2);
  \node[boxlab, anchor=north west, align=center] at (3.35,2.05) {string-level\\aux.\ columns};

  %==================== Panel 2: char-level columns (placed to the right) ====================
  \begin{scope}[xshift=8.1*\u cm, yshift=5.2*\u cm]
  % group (string) labels above the char row
  \foreach \cx/\nm/\col in {2/apple/applecol, 7.5/orange/orangecol, 13.5/cherry/cherrycol, 19.5/banana/bananacol}{
    \node[font=\ttfamily, text=\col] at (\cx*0.6, 1.55) {\nm};
  }

  % row labels
  \node[rowlab] at (-0.5, 0.00){$\mathsf{char}$};
  \node[rowlab] at (-0.5,-0.68){$\mathsf{orig\text{-}ind}$};
  \node[rowlab] at (-0.5,-1.36){$\mathsf{int\text{-}ind}$};
  \node[rowlab] at (-0.5,-2.04){$\mathsf{char\text{-}act}$};
  \node[rowlab] at (-0.5,-2.72){$\mathsf{bnd}$};

  % position indices (set smaller than the body text)
  \foreach \x in {0,...,22}{\node[idx] at (\x*0.6,0.85) {\x};}

  % pos / char / orig / o / a_c / s_b ; cherry (orig 2) is inactive so a_c=0 there.
  \foreach \x/\ch/\og/\oo/\ac/\sb in {
     0/a/0/0/1/1,  1/p/0/1/1/0,  2/p/0/2/1/0,  3/l/0/3/1/0,  4/e/0/4/1/0,
     5/o/1/0/1/1,  6/r/1/1/1/0,  7/a/1/2/1/0,  8/n/1/3/1/0,  9/g/1/4/1/0, 10/e/1/5/1/0,
    11/c/2/0/0/1, 12/h/2/1/0/0, 13/e/2/2/0/0, 14/r/2/3/0/0, 15/r/2/4/0/0, 16/y/2/5/0/0,
    17/b/3/0/1/1, 18/a/3/1/1/0, 19/n/3/2/1/0, 20/a/3/3/1/0, 21/n/3/4/1/0, 22/a/3/5/1/0}{
    \pgfmathsetmacro{\px}{\x*0.6}
    \edef\gc{\ifcase\og applecol\or orangecol\or cherrycol\or bananacol\fi}
    \node[cell, fill=\gc!15] at (\px, 0.00) {\texttt{\ch}};
    \node[cell, fill=\gc!15] at (\px,-0.68) {\og};
    \node[cell, fill=\gc!15] at (\px,-1.36) {\oo};
    \node[cell, fill=\gc!15] at (\px,-2.04) {\ac};
    \node[cell, fill=\gc!15] at (\px,-2.72) {\sb};
  }

  % padding tail: one padding slot, an ellipsis, then the final index
  % (31 = last cell of the size-32 hypercube that M=23 pads up to).
  \node[idx] at (23*0.6,0.85) {23};
  \foreach \yy in {0.00,-0.68,-1.36,-2.04,-2.72}{
    \node[cell] at (23*0.6,\yy) {0};
  }
  \foreach \yy in {0.00,-0.68,-1.36,-2.04,-2.72}{
    \node[font=\scriptsize] at (24*0.6,\yy) {$\cdots$};
  }
  \node[idx] at (25*0.6,0.85) {31};
  \foreach \yy in {0.00,-0.68,-1.36,-2.04,-2.72}{
    \node[cell] at (25*0.6,\yy) {0};
  }

  % dotted box around the char-level auxiliary columns
  \draw[black, thick, densely dotted, rounded corners=2pt] (-2.45,0.50) rectangle (15.4,-3.19);
  \node[boxlab] at (6.25,-3.34) {character-level auxiliary columns};
  \end{scope}

\end{tikzpicture}
\caption{String arithmetization of a four-row column
$\strcol=(\texttt{apple},\texttt{orange},\texttt{cherry},\texttt{banana})$ with
\texttt{cherry} inactive.}
\label{fig:string-arith-example}
\ifsp\end{figure*}\else\end{figure}\fi



As shown in \cref{fig:string-arith-example}, we keep the hash column $(\act,\mathsf{h})$ of $\truthtable$ as \emph{the main string column}, and view it as the $\strcol.\mathsf{main}$. To this we add \emph{auxiliary string columns}, which come in two types with two different lengths:

\parhead{Type 1: String-level auxiliary columns} These are columns of length $N$---the length of the main column, with one entry per string. Together with the activator $\mathsf{a}$ and the index column $\mathsf{ind}$ of the main table, which they share, they form the \emph{auxiliary string table} $\strcol.\mathsf{aux\_strs}$:
\begin{itemize}[nosep]
    \item \emph{Length column $\mathsf{l}$}, with $\mathsf{l}[i]$ being the length of the $i$-th string in $\strcol$, counted in bytes of its \textsc{utf-8} encoding.
    \item \emph{Rank column $\mathsf{rk}$}, with $\mathsf{rk}[i]$ being the position of the $i$-th string in the lexicographic order of the distinct strings of $\strcol$, so that equal strings share a rank and comparing ranks compares strings. We use it in \cref{sec:rank} to sort strings at the cost of sorting integers.
    \end{itemize}
    \parhead{Type 2: Character-level auxiliary columns} These are columns of length $M := \sum_{s \in \strcol} |s|$, the total number of characters in the column. They form the \emph{auxiliary character table} $\strcol.\mathsf{aux\_chars}$, with one row per character:
   \begin{itemize}[nosep]
    \item \emph{Character column $\mathsf{char}$} contains the individual characters of all strings in $\strcol$ concatenated to each other. Throughout, a \emph{character} is one byte of the string's \textsc{utf-8} encoding, so every entry of $\mathsf{char}$ lies in the alphabet $\Sigma = \{0, \dots, 2^{\alpha}-1\}$ with $\alpha = 8$, and a multi-byte code point occupies several consecutive entries.
    \item \emph{Origin index column $\mathsf{orig\text{-}ind}$}, with $\mathsf{orig\text{-}ind}[j]$ being the index of the string in $\strcol$ to which the $j$-th character of $\mathsf{char}$ belongs.
    \item \emph{Internal index column $\mathsf{int\text{-}ind}$}, with $\mathsf{int\text{-}ind}[j]$ being the position of the $j$-th character of $\mathsf{char}$ within its own string, counting from $0$. Equivalently, $\mathsf{int\text{-}ind}$ resets to $0$ at every string boundary and increments by one across the characters of each string.
    \item \emph{Character activator column $\mathsf{char\text{-}act}$}, with $\mathsf{char\text{-}act}[j]$ being $1$ if the $j$-th character of $\mathsf{char}$ is active, and $0$ otherwise.
    \item \emph{Boundary marker column $\mathsf{bnd}$}, with $\mathsf{bnd}[j]$ being $1$ if the $j$-th character of $\mathsf{char}$ is the first character of a string (a string boundary), and $0$ otherwise.
    \end{itemize}
\cref{fig:string-arith-example} illustrates this arithmetization on a sample four-row string-typed column. Our protocols treat the three tables as ordinary relations.

\parhead{Abuse of notation} Since $\strcol.\mathsf{main}$ and $\strcol.\mathsf{aux\_strs}$ share their rows, we let a projection (\cref{sec:prelims}) of $\strcol$ itself, such as $\strcol[\mathsf{h}, \mathsf{l}, \mathsf{rk}]$, draw its columns from both.

\parhead{Redundant columns} Some of the columns above are completely redundant: they add nothing to the characterization of a string, which $\mathsf{char}$, $\mathsf{orig\text{-}ind}$ and $\mathsf{int\text{-}ind}$ already determine uniquely. Including them is purely a form of preprocessing that speeds up the prover in the protocols of the following sections. The boundary marker $\mathsf{bnd}$, for instance, is just the indicator of $\mathsf{int\text{-}ind} = 0$, which the prover could compute and send on the fly; but since it would then have to prove that the indicator is correct in every protocol that uses it, we let the data owner ship it once as part of the auxiliary columns instead. The same holds for the length $\mathsf{l}$, which is determined by the character rows of each string, and for the rank $\mathsf{rk}$, which is determined by the strings themselves.
